\documentclass[11pt,a4paper]{article}
\usepackage[T1]{fontenc}
\usepackage[margin=24mm]{geometry}
\usepackage{amsmath,amssymb,booktabs,longtable,array,enumitem,microtype}
\usepackage{xcolor,xurl}
\usepackage[colorlinks=true,linkcolor=blue!45!black,urlcolor=blue!45!black]{hyperref}
\newcommand{\code}[1]{\ifmmode\mbox{\ttfamily\detokenize{#1}}\else\nolinkurl{#1}\fi}
\newcommand{\src}[1]{\hyperlink{src:#1}{\textsf{[#1]}}}
\newcommand{\GeV}{\,\mathrm{GeV}}
\newcommand{\True}{\mathsf{True}}
\newcommand{\False}{\mathsf{False}}
\newcommand{\Unknown}{\mathsf{None}}
\newcommand{\BR}{\operatorname{BR}}
\newcommand{\dd}{\mathrm{d}}
\setlist{itemsep=3pt,topsep=5pt}
\setlength{\parskip}{4pt}
\setlength{\parindent}{0pt}
\setlength{\emergencystretch}{2em}
\renewcommand{\arraystretch}{1.14}
\title{Constraints and selection logic\\for the two-singlet dark-matter scan}
\author{Repository implementation note}
\date{1 October 2026}

\begin{document}
\maketitle
\begin{abstract}
We document the constraints implemented in the active $v_x=0$ scan and their
use in the comprehensive plotting suite. We distinguish the order of numerical
evaluation, the logical definition of viability, and the cumulative sequence
shown in the plots. We also specify the independent vacuum, running-coupling
and finite-temperature assessments, together with the treatment of missing
results and the additional restrictions used to display particular observables.
The definitions below describe the checked-out source, rather than a new
choice of physics cuts.
\end{abstract}

\textbf{Scope and provenance.}
The repository is \code{twosingletDM}, inside the parent
\code{TwoSingletDM} directory. The active prescription is
\code{trsm_constraints_v2}, schema 2. This note was prepared from the working
tree on the date above, whose base Git commit was
\code{c29d1c28d4d39f84168499c8feaadc7a0cd75511}; local modifications were present,
so the commit alone does not identify every inspected file. Source references
such as \src{P} resolve to the file/function index in
Section~\ref{sec:sources}. The main reference for plot selections is
\code{plot_trsm_constraint_suite.py}, including its companion
\code{plot_trsm_v2.py}. The general $v_x\ne0$ model and historical plotters
remain in the repository, but do not define the active dark-matter selection.
No particular scan file was selected for this note; survival counts therefore
depend on the input dataset and are not quoted.

\tableofcontents

\section{The three meanings of constraint order}
\label{sec:order}

\subsection{Numerical evaluation in the current generator}
For a normally evaluated point,
\code{_evaluate_trsm_point_vxzero} performs the following operations.
The calculations are recorded independently: an early exclusion does not
normally prevent the remaining zero-temperature checks from being evaluated.
\src{G}
\begin{enumerate}
\item Construct masses, mixing factors, scalar couplings, partial widths,
branching fractions and the physical invisible widths.
\item Evaluate the correlated $S,T,U$ test, then the tabulated $W$-mass test.
\item Evaluate HiggsBounds and HiggsSignals through HiggsTools.
\item Evaluate the tree-level boundedness and unitarity test
(\code{thc}), then the tree-vacuum and RG diagnostics.
An explicitly unbounded vacuum status also sets \code{thc=False}.
The stored \code{evo} value is the RG integration-success flag.
\item Run the dark-matter calculation and assess relic density, direct
detection, gamma-ray lines and, when enabled, CMB energy injection.
\item Evaluate the independent flavour constraint. Construct the experimental,
DM and EWPT-eligibility subsets, and determine non-DM and full viability.
\item If requested and eligible, calculate the selected MadGraph rates.
\item If requested and EWPT-eligible, run BSMPT and append its candidate,
phase-history and freeze-out diagnostics.
\item Write the point to the main scan ledger, including excluded and
unassessed points. An optional DM-failed sidecar contains points with
\code{dm=False} that pass all non-DM constraints.
\end{enumerate}
An input/calculation exception can instead produce an unassessed row with
its reason. The complete-ledger policy does not mean that every calculation
succeeded or that every written point is viable. The historical
\code{--nrandom-count-evo-thc} option changes the stopping counter to
\code{evo=True} and \code{thc=True}; it does not redefine full viability.

\subsection{Logical selections: what actually has to pass}
\label{sec:logical}
Define $T$ to mean an explicit tree-theory pass, $E$ the assessed experimental
pass, $F$ the flavour pass and $D$ the aggregate DM pass. In v2,
\begin{align}
 T &\equiv (\code{thc}\text{ is }\True),\\
 E &\equiv (\code{experimental_subset}\text{ is }\True),\\
 N &\equiv T\land E\land F,
 & V &\equiv T\land E\land F\land D.
 \label{eq:viability}
\end{align}
Here $N$ is \code{non_dm_viability} and $V$ is \code{full_viability} in the
plotting suite. The experimental subset combines HiggsBounds, HiggsSignals
and $S,T,U$, with the $W$-mass comparison added where its table covers $M_2$.
Flavour is a separate factor; it is not part of
\code{experimental_subset}. \code{dm_subset} copies the aggregate
\code{dm} verdict. \src{G}\src{R}\src{P}

The separate EWPT entry gate is
\begin{equation}
 \code{ewpt_eligible}
 = (\code{thc}\text{ is }\True)
 \land (\code{experimental_subset}\text{ is not }\False).
 \label{eq:ewptgate}
\end{equation}
Thus an incomplete experimental assessment is allowed into EWPT exploration.
DM, flavour, a nonglobal tree vacuum, RG failure and the strict-theory subset
do not veto that exploration. Actual execution also requires the relevant
EWPT run option. Full viability requires an explicit experimental pass,
whereas EWPT eligibility admits an unassessed experimental result.

\begin{longtable}{@{}p{0.25\textwidth}p{0.36\textwidth}p{0.33\textwidth}@{}}
\toprule
Assessment & Definition or role & Required for current full viability?\\
\midrule\endhead
Tree theory & BFB and perturbative unitarity; explicitly unbounded directions fail & Yes\\
Experimental & HB, HS, STU, and covered $W$ mass & Yes\\
Flavour & Covered $\Upsilon$ lepton-product limits & Yes; non-veto outside coverage\\
Dark matter & Relic, DD, line and enabled CMB prescription & Yes\\
Global tree vacuum & Desired vacuum globally minimal & Separate subset\\
Running couplings & Integration, BFB and unitarity up to $1$ TeV & Separate subset\\
EWPT candidates & Qualitative baryogenesis/GW criteria & Separate flags\\
Thermal DM checks & Equilibrium $Z_2$ compatibility and VEV/coupling variation & Diagnostics\\
Widths and lifetimes & Physical $\Gamma_i/M_i$ and $c\tau_i$ & Diagnostics\\
\bottomrule
\end{longtable}

MadGraph is run, when requested, on $V$ by default. The
\code{--mg5-without-dm} option replaces this by $N$, dropping only DM.
The active v2 call site implements this directly; the older helper
\code{mg5_point_eligible} still has an \code{evo} argument and is not the
definition used by that call site. \src{G}

\subsection{The cumulative order displayed in the plotting suite}
\label{sec:cumulative}
The cumulative mass/parameter plots deliberately use the following sequence:
\begin{align}
 C_0 &= \{\text{all stored rows}\},\nonumber\\
 C_1 &= C_0\cap\{\code{hb=True}\},\nonumber\\
 C_2 &= C_1\cap\{\code{hs=True}\},\nonumber\\
 C_3 &= C_2\cap\{\code{wmass=True}\},\nonumber\\
 C_4 &= C_3\cap\{\code{dm=True}\}.
 \label{eq:plotcutflow}
\end{align}
This is implemented in \code{cumulative_constraint_masks}. It supplies plots
24--29, \code{dashboard_cumulative_constraint_summary}, and the
\code{cumulative_selection_*} rows in \code{constraint_summary.tsv}.
\textbf{The cumulative sequence omits EWPO, tree theory and flavour.}
Consequently, $C_4$ is not the full-viability set $V$. It also does not
represent the generator's execution order. Plot 22,
\code{22_constraint_counts}, reports independent totals, rather than a
sequential cutflow. \src{P}

\subsection{Pass, exclusion and unavailable results}
\label{sec:nullable}
The assessment layer uses three values: \code{True}, \code{False}, and
\code{None} (\code{null} in JSON and usually \code{nan} in TSV).
Its nullable conjunction returns false if any input is explicitly false,
true if every input is explicitly true, and unavailable otherwise.
Thus a known exclusion survives another component's missing assessment.
\src{I}

The main plotter converts explicit true values into pass masks. It keeps
availability masks for detailed nullable columns, but its simple binary
status views can group unassessed results with nonpassing results. A
nonpass count therefore need not be an exclusion count. The dedicated v2
candidate panels retain true/false/unknown distinctions. Missing flavour
results cannot satisfy $N$ or $V$; the separate
\code{pre_flavour_full_viability} mask is $T\land E\land D$ and only
reproduces the selection before flavour was added. \src{P}\src{V}

\section{Model conventions and scan restrictions}
\label{sec:model}
We consider the $v_x=0$ branch with $a_{13}=a_{23}=0$.
The stable scalar is $h_3=X$, with mass $M_3$ and vanishing doublet projection.
For $\theta=a_{12}$, $c=\cos\theta$ and $s=\sin\theta$,
\begin{equation}
 h_1=c\,\phi-s\,s_0,\qquad h_2=s\,\phi+c\,s_0,
 \qquad k_1=c,\quad k_2=s,\quad k_3=0,
\end{equation}
where $\phi,s_0$ denote fluctuations about $v,v_s$.
The shared inputs are
\begin{equation}
 M_1=125.09\GeV,\quad
 G_F=1.1663787\times10^{-5}\GeV^{-2},\quad
 M_W=80.385\GeV,\quad M_Z=91.1876\GeV,
\end{equation}
with $v=(\sqrt{2}G_F)^{-1/2}=246.2196508\GeV$.
These are the repository's fixed inputs, not a claim about the most recent
experimental averages. \src{I}\src{W}

The unshifted real-field potential convention, with $\varphi=(h,s,x)$, is
\begin{equation}
 V(\varphi)=\frac12\sum_i\mu_i^2\varphi_i^2
 +\frac14\sum_i\lambda_i\varphi_i^4
 +\frac14\sum_{i<j}\lambda_{ij}\varphi_i^2\varphi_j^2.
 \label{eq:potential}
\end{equation}
The reconstructed quartics are
\begin{align}
 \lambda_\Phi&=\frac{M_1^2c^2+M_2^2s^2}{2v^2},&
 \lambda_S&=\frac{M_1^2s^2+M_2^2c^2}{2v_s^2},&
 \lambda_{\Phi S}&=\frac{(M_2^2-M_1^2)sc}{vv_s}.
\end{align}
The independent inputs are $M_2,M_3,v_s,\theta,\lambda_X,
\lambda_{\Phi X},\lambda_{SX}$. In particular,
\begin{align}
 \mu_\Phi^2&=-\lambda_\Phi v^2-\tfrac12\lambda_{\Phi S}v_s^2,\nonumber\\
 \mu_S^2&=-\lambda_S v_s^2-\tfrac12\lambda_{\Phi S}v^2,\nonumber\\
 \mu_X^2&=M_3^2-\tfrac12\lambda_{\Phi X}v^2-\tfrac12\lambda_{SX}v_s^2.
\end{align}
\src{T}\src{Q}

\subsection{Sampling priors are separate from exclusions}
The random-scan defaults are $M_2,M_3\in[4,1000]\GeV$,
$v_s\in[50,500]\GeV$, $k_1\in[0.90,1]$ and
$\lambda_X\in[0,10]$. Both signs of
$\theta=\pm\arccos k_1$ are sampled. Each default portal is drawn as
$u\,z$, with $u$ uniform on $[0,0.002]$ and $z$ uniform on $[-1,1]$.
Its support is consequently $[-0.002,0.002]$, with a nonuniform density;
the configured positive magnitude interval does not prohibit negative portals.
\src{G}

The generator's ordinary mass mode draws
\begin{align}
 M_2&\in[m_{2,\min},\min(m_{2,\max},m_{3,\max}-M_1)],\nonumber\\
 M_3&\in[\max(m_{3,\min},M_2+M_1),m_{3,\max}],
\end{align}
so that $M_3\ge M_2+M_1$. This is a sampling condition.
The following alternatives replace it:
\begin{itemize}
\item \code{--independent-m3}: independent rectangular mass sampling.
\item \code{--resonantDM1} or \code{--resonantDM2}: respectively
$2M_3=M_1$ or $2M_3=M_2$, subject to the configured ranges.
\item \code{--approximate-resonantDM}: sample one of the available
$|2M_3-M_1|\le\delta_{\rm res}$ and
$|2M_3-M_2|\le\delta_{\rm res}$ branches, with equal branch probability
when both exist and with range-compatible conditional sampling.
\item \code{--mass-ratio-m3-2m2}: $|M_3-2M_2|\le\delta_{\rm res}$.
This is a separate mass-ratio study, not the annihilation-pole relation.
\end{itemize}
All mass and coupling ranges can be overridden. The supplied
\code{config/next-scan-v2.json} chooses approximate annihilation resonance
sampling with $\delta_{\rm res}=2\GeV$, the ranges above,
micrOMEGAs 7.1.4 with CMB enabled, and an EWPT upper temperature of
$1000\GeV$. The generic EWPT command-line default is instead $300\GeV$.
\src{G}\src{C}

The optional \code{--scan-k133-k233} mode samples each $K$ as a uniform
positive magnitude times a uniform number in $[-1,1]$, with magnitude
intervals $[10^{-6},8]\GeV$ and $[10^{-4},8]\GeV$.
\code{--scan-k133-k233-log} uses independent random signs and uniform
$\log_{10}(|K|/\mathrm{GeV})$ in $[-4,3]$ and $[-3,5]$.
Both modes convert to portal inputs using the inverse of
Eq.~\eqref{eq:portals}; they are alternative sampling priors.

\section{Tree theory, global vacuum and running couplings}
\subsection{Boundedness from below and tree-level unitarity}
Introduce the symmetric matrix $A$, with diagonal
$(\lambda_\Phi,\lambda_S,\lambda_X)$ and off-diagonal entries
$A_{ij}=\lambda_{ij}/2$, and
$\overline A_{ij}=A_{ij}+\sqrt{A_{ii}A_{jj}}$.
The quartic boundedness test imposes copositivity:
\begin{gather}
 A_{ii}\ge0,\qquad \overline A_{12},\overline A_{13},\overline A_{23}\ge0,\\
 \sqrt{A_{11}A_{22}A_{33}}
 +A_{12}\sqrt{A_{33}}+A_{13}\sqrt{A_{22}}+A_{23}\sqrt{A_{11}}
 +\sqrt{2\overline A_{12}\overline A_{13}\overline A_{23}}\ge0.
\end{gather}
Each comparison allows a numerical tolerance of $10^{-12}$ below zero;
diagonal values within that tolerance are clamped to zero in the radicals.
\src{T}

The strict perturbative bounds are
\begin{equation}
 |\lambda_\Phi|,|\lambda_S|,|\lambda_X|<4\pi,\qquad
 |\lambda_{\Phi S}|,|\lambda_{\Phi X}|,|\lambda_{SX}|<8\pi,
\end{equation}
and every eigenvalue $u_a$ of
\begin{equation}
 U=\begin{pmatrix}
 12\lambda_\Phi&2\lambda_{\Phi S}&2\lambda_{\Phi X}\\
 2\lambda_{\Phi S}&6\lambda_S&\lambda_{SX}\\
 2\lambda_{\Phi X}&\lambda_{SX}&6\lambda_X
 \end{pmatrix}
\end{equation}
must satisfy $|u_a|<16\pi$. Equality fails these strict bounds.
Finite inputs and nonzero $v_s$ are required. The real symmetric eigenvalue
routine \code{eigvalsh} is used. Their conjunction supplies \code{thc},
with the unbounded-direction correction described below. \src{T}\src{G}

\subsection{The desired vacuum}
The tree-vacuum diagnostic compares the desired point
$(h,s,x)=(v,v_s,0)$ with stationary points on all eight supports:
the origin, three axes, three planes and the full three-field branch.
Writing $y_i=\varphi_i^2\ge0$, it solves
$A_{\mathcal S}y_{\mathcal S}=-\mu_{\mathcal S}^2$ on each support,
checks the Hessian and compares the potential depths of local minima.
The target must be local and obey
\begin{equation}
 V_{\rm target}-V_{\min}\le
 10^{-8}\max(1,|V_{\rm target}|,|V_{\min}|),
\end{equation}
where the numerical potential values are in $\mathrm{GeV}^4$.
Degenerate global minima pass. Singular systems are examined for feasible
flat stationary families and for nonnegative quartic-null directions with
a negative quadratic contribution. An unresolved flat family remains
unassessed. The full implementation includes coordinate and Hessian
tolerances. \src{Q}

\code{vacuum_tree_global} is an independent subset. A known deeper minimum
does not by itself fail \code{thc}. However, the statuses
\code{unbounded_quartic} and \code{unbounded_flat_direction} do force
\code{thc=False}. The BSMPT NLO stability status is stored separately from
this tree-level result. \src{G}\src{Q}

\subsection{RG screening through 1 TeV}
One-loop running starts at $\mu_0=91\GeV$ and extends to
$\mu_{\max}=1000\GeV$. The boundary values are
\begin{equation}
 g_3=1.221,\qquad g_2=\sqrt{0.424},\qquad
 g_1=\sqrt{0.1273},\qquad y_t=0.96738,
\end{equation}
with the scalar couplings reconstructed from tree masses and mixing.
There is no threshold matching or precision pole-to-running matching.
The integrator uses relative/absolute tolerances $10^{-8}/10^{-10}$ and
maximum step $(1000-91)/256\GeV$. BFB and the same strict unitarity tests
are checked on solver nodes and 513 logarithmically spaced samples.
The first detected violation is refined to a bracket narrower than
$10^{-5}\GeV$. This is a numerical trajectory screen. \src{I}\src{Q}

\code{rg_integration_success}, \code{rg_bfb} and \code{rg_unitarity}
are distinct. An observed physical violation is false even if integration
later fails; with no detected violation, a failed integration leaves the
corresponding physical verdict unavailable. The strict subset
\code{theory_strict_subset} is the nullable conjunction of
\code{vacuum_tree_global}, \code{rg_integration_success},
\code{rg_bfb} and \code{rg_unitarity}.
It is not required by Eq.~\eqref{eq:viability}. In the current generator,
\code{evo} records integration success; it is not a synonym for either
running BFB or running unitarity.

\section{Higgs, electroweak and flavour constraints}
\subsection{Physical widths and HiggsTools inputs}
The portal coefficients in the scalar potential are
\begin{align}
 K_{133}&=\tfrac12(\lambda_{\Phi X}vc-\lambda_{SX}v_ss),\nonumber\\
 K_{233}&=\tfrac12(\lambda_{\Phi X}vs+\lambda_{SX}v_sc).
 \label{eq:portals}
\end{align}
Here $K_{i33}$ multiplies $h_i h_3^2$, so the identical-scalar vertex has
magnitude $2K_{i33}$ and
\begin{equation}
 \Gamma(h_i\to h_3h_3)=\frac{K_{i33}^2}{8\pi M_i}
 \sqrt{1-\frac{4M_3^2}{M_i^2}}\quad(M_i>2M_3).
\end{equation}
The width is zero at and below threshold. Invisible and open scalar-cascade
widths enter the physical total widths and dilute visible branching
fractions. The same convention applies to other identical-scalar decays.
\src{W}

HiggsBounds receives the mass, rescaled SM couplings and decay information
and supplies the stored \code{hb} verdict through its \code{allowed}
result. The relevant analyses are those in the installed, fingerprinted
HiggsBounds dataset; there is no repository-wide single cross-section cut
that replaces this provider result. \src{H}

HiggsSignals supplies
\begin{equation}
 \Delta\chi^2_{\rm HS}=\chi^2_{\rm point}-\chi^2_{\rm SM}<4.00.
 \label{eq:hs}
\end{equation}
The SM reference is constructed with the same mass and SM decay-input
pipeline. Nonfinite values are unassessed. This strict, fixed screening
threshold is not a profiled multi-parameter confidence region.
The active path includes physical invisible decays in HiggsTools; it does
not add an independent fixed upper cut on $\BR(h_1\to\mathrm{invisible})$.
A provider width floor of $10^{-10}\GeV$ is numerical and is not a
physical acceptance cut. \src{H}

\subsection{Correlated electroweak precision observables}
The active fit allows nonzero $U$. The singlet shifts have the form
\begin{equation}
 \Delta O=s^2[O(M_2)-O(M_1)],\qquad O\in\{S,T,U\}.
\end{equation}
The code evaluates
\begin{equation}
 \chi^2_{STU}=(\boldsymbol\Delta-\boldsymbol c)^T
 \Sigma^{-1}(\boldsymbol\Delta-\boldsymbol c)\le7.82,
 \qquad \Sigma_{ij}=\rho_{ij}\sigma_i\sigma_j,
\end{equation}
using, in $(S,T,U)$ order,
\begin{align}
 \boldsymbol c&=(-0.045235,\;0.013373,\;-0.016131),\nonumber\\
 \boldsymbol\sigma&=(0.094703,\;0.11518,\;0.087103),\nonumber\\
 \rho&=\begin{pmatrix}1&0.048&0.747\\0.048&1&-0.162\\0.747&-0.162&1\end{pmatrix}.
\end{align}
The repository attributes these inputs to a Jens Erler private communication
dated 16 May 2025. They are documented here as implemented inputs, rather
than an independently authenticated public likelihood. The loop normalization
retains $\alpha(0)=7.2973525693\times10^{-3}$. A nonfinite test is
unassessed. Historical two-dimensional $U=0$ and projected fits in the
file are not called by the active generator. \src{E}

\subsection{The separate W-mass curve}
For $133\le M_2/\mathrm{GeV}\le999$, the code imposes
\begin{equation}
 |\sin\theta|\le s_{\max}^{W}(M_2),
\end{equation}
where $s_{\max}^{W}$ is the existing cubic interpolation of
\code{datafiles/Tania_MW_SnowmassWhitepaper.dat}. Equality passes.
There is no extrapolation. Finite points outside the table range receive
a true non-veto from the checker and
\code{wmass_status=outside_coverage}; the v2 experimental subset omits
the $W$-mass factor there. Invalid inputs are unassessed. The STU and
$W$-mass prescriptions are applied as screens, without forming a joint
likelihood or asserting their statistical independence. \src{E}\src{R}

\subsection{Upsilon lepton-product bounds}
The independent \code{upsilon_leptons_v1} prescription tests
\begin{equation}
 P_{i\ell}=\BR(\Upsilon(1S)\to\gamma h_i)
             \BR(h_i\to\ell^+\ell^-),\qquad \ell=\mu,\tau.
\end{equation}
The implemented production formula is
\begin{align}
 \frac{\BR(\Upsilon\to\gamma h_i)}{\BR(\Upsilon\to e^+e^-)}
 &=\frac{G_Fm_b^2}{\sqrt2\pi\alpha}\,|k_i|^2
 \left(1-\frac{M_i^2}{m_\Upsilon^2}\right)F_{bS}(M_i),\nonumber\\
 F_{bS}(M_i)&=\frac23\left[1-\left(\frac{M_i}{m_\Upsilon}\right)^6\right],
 \qquad M_i<m_\Upsilon.
\end{align}
It is zero when on-shell production is forbidden. The retained central
inputs are $m_b=4.18\GeV$, $m_\Upsilon=9.46040\GeV$,
$\alpha^{-1}=137.035999084$ and
$\BR(\Upsilon\to e^+e^-)=0.0239$. Visible lepton BRs are multiplied by
$\Gamma_{\rm base}/\Gamma_{\rm physical}$, so invisible decays dilute the
signal without altering the underlying base BR arrays. \src{F}

For each lepton channel, the strongest available individual digitized
Belle or BaBar 90\% CL limit is used. The nominal windows are inclusive
$[0.212,9.2]\GeV$ for muons and $[3.55,9.2]\GeV$ for taus, additionally
restricted to each CSV curve's coverage. The supplied tau curves begin at
approximately $3.834$ and $3.839\GeV$. Interpolation is linear in mass
and logarithmic in the limit, with no extrapolation. Every covered ratio
$P_{i\ell}/L_\ell(M_i)$ must be at most one; equality passes.
These are the retained prompt, narrow-resonance bounds, not a combined
experimental likelihood. No extra lifetime or broad-width acceptance
test is imposed. \src{F}

Outside coverage and exactly zero production are non-vetoes. A missing
required decay input is unavailable; in particular, the base decay provider
starts at $4\GeV$. In the active branch, $h_1$ is above the search range
and $h_3$ has zero doublet projection, leaving a light $h_2$ as the relevant
decaying scalar. The nullable conjunction over the scalar assessments
defines \code{flavour}, and \code{flavour_max_ratio} records the largest
covered ratio.

\section{Dark-matter constraints}
\label{sec:dm}
The default backend is micrOMEGAs 7.1.4; 6.1.15 is an explicit comparison
backend. A usable result requires a successful solver and finite physical
outputs. The aggregate DM verdict is unavailable on a solver error, even
if some raw diagnostics were printed. The active calculation uses the
abundance fraction
\begin{equation}
 \xi=\min\left(1,\frac{\Omega h^2}{0.12}\right).
 \label{eq:xi}
\end{equation}
Direct detection is rescaled linearly and annihilation observables
quadratically. The abundance reference $0.12$ is distinct from the relic
upper bound $0.121$. \src{D}\src{I}

\subsection{Relic abundance}
The required upper-bound condition is
\begin{equation}
 0\le\Omega h^2\le0.121.
\end{equation}
There is no positive lower abundance cut and no requirement that $h_3$
provide all of the dark matter. Negative/nonfinite solver output is invalid,
rather than an accepted underabundant point. The nominal freeze-out
temperature is recorded as $T_f=M_3/X_f$ when a valid $X_f>0$ is available.
\src{D}

\subsection{Spin-independent direct detection}
The default is the stored LZ WS2024 observed, power-constrained 90\% CL
elastic isoscalar per-nucleon table. The wrapper reads the neutron SI cross
section from the micrOMEGAs cross-section output, in pb. The stored
experimental values in $\mathrm{cm}^2$ are converted using
$1\,\mathrm{cm}^2=10^{36}\,\mathrm{pb}$. The cut is
\begin{equation}
 \xi\,\sigma_n^{\rm SI}\le L_{\rm LZ}(M_3),
 \quad\text{equivalently}\quad
 \sigma_n^{\rm SI}\le L_{\rm LZ}(M_3)/\xi\quad(\xi>0).
\end{equation}
The 26-point default table covers $9$--$10000\GeV$ and is interpolated
logarithmically in both mass and limit. There is no extrapolation.
For positive $\xi$, an uncovered DD assessment is unavailable and cannot
supply a DM pass. At $\xi=0$, the implementation sets the comparison limit
to infinity, so a finite cross section is a non-veto. \src{D}\src{L}

\code{dm_dir_det} stores the raw cross section, while
\code{dm_dir_det_limit} stores the abundance-rescaled limit.
The plotting ratio is their quotient; it is not a comparison of an already
rescaled cross section with a second rescaled bound. An explicit
\code{--dm-limit-table} can replace the default data. The optional LZ2026
high-mass approximation and historical fitted providers are not the default
WS2024 prescription and must be identified from the saved metadata.

\subsection{Gamma-ray lines}
The retained indirect-detection screen compares the individual line fluxes
with the tabulated Fermi-LAT R16 line limit:
\begin{equation}
 R_{\rm line}=\max_{a\,\mathrm{usable}}
 \frac{\xi^2\Phi_a}{L_{\rm R16}(E_{\gamma,a})}\le1.
\end{equation}
The energy coverage is $0.214$--$462\GeV$, with log--log interpolation
and no extrapolation. The driver supplies photon-containing two-body
annihilation channels, with $E_\gamma=M_3$ for $\gamma\gamma$ and
$E_\gamma=M_3(1-m_Y^2/(4M_3^2))$ for $\gamma Y$.
The implementation takes the largest individual line ratio; it does not
sum nearby line fluxes into a new combined bound. The driver approximates
the R16 region using $0.5^\circ$ patches inside
$\sqrt{l^2+b^2}\le16^\circ$, excluding $|b|<5^\circ$ with
$|l|>6^\circ$, and weights the flux by photon multiplicity.
This is the retained line-flux prescription, not a published-ROI likelihood
evaluation. \src{D}

If no usable line is found, the code stores
\code{dm_indirect_available=False} and
\code{dm_indirect_detection_excluded=False}. This is a non-veto in the
aggregate DM verdict, unlike an unavailable covered DD or enabled CMB
assessment. The line availability must therefore be reported alongside
the aggregate flag. The printed continuum flux is a diagnostic and is not
an additional bound in this selection. \src{D}

\subsection{CMB energy injection}
New version-7 scans enable the Planck CMB screen by default; version-6
comparison runs and legacy resumes without saved CMB settings default to
disabled. The explicit switches are \code{--planck-cmb} and
\code{--no-planck-cmb}. When enabled,
\begin{equation}
 R_{\rm CMB}=\xi^2 R_{\rm CMB}^{\rm raw}\le1,
 \qquad
 p_{\rm ann}^{\max}=3.2\times10^{-28}
 \,\mathrm{cm}^3\,\mathrm{s}^{-1}\,\mathrm{GeV}^{-1}.
\end{equation}
The driver uses the built-in Planck2018 spectral energy-deposition
prescription in a low-velocity $s$-wave approximation, with the quoted
95\% CL normalization. A missing or invalid enabled CMB calculation is
unassessed; it cannot supply a pass. A disabled CMB factor is omitted.
\src{B}

\subsection{Aggregate DM verdict and equality conventions}
For a successful solver, \code{DMSummary.passed} forms a nullable
conjunction of the relic pass, DD pass, absence of a line exclusion, and
the CMB pass when enabled. Exclusions use strict ``greater than'' comparisons,
so equality at the relic, DD, line or CMB upper limit passes.
A known failing component gives \code{dm=False}; otherwise missing required
components give \code{dm=None}. Valid component information is retained.
The plotting suite trusts this stored aggregate flag and does not rebuild
it from the displayed ratios. \src{D}\src{P}

\section{EWPT entry, candidate criteria and thermal diagnostics}
\subsection{Optional quartic preselection before BSMPT}
\code{--ewpt-require-eq418} is off by default. If enabled in the active
EWPT path, the repository's ``Eq.~4.18'' test must pass before BSMPT is
called. Define the matrix $C$ with diagonal
$c_i=\lambda_i$ and off-diagonal $c_{ij}=\lambda_{ij}/2$.
The implemented conditions are
\begin{equation}
 c_i>0,\qquad c_ic_j-c_{ij}^2>0\quad(i<j),\qquad \det C>0,
\end{equation}
where
\begin{equation}
 \det C=c_\Phi c_S c_X+2c_{\Phi S}c_{\Phi X}c_{SX}
 -c_{\Phi S}^2c_X-c_{\Phi X}^2c_S-c_{SX}^2c_\Phi.
\end{equation}
All inequalities are strict. This optional positive-definiteness screen
uses the same matrix normalization as $A$ in Section~3, but imposes a stronger
condition than copositivity. The diagonal normalization was corrected on
29 October 2026; earlier code used $c_i=\lambda_i/2$, which can change the
sign of $D$ and the preselection decision. Existing plots and stored outcomes
must be recomputed to use the corrected normalization.
The screen gates EWPT execution and does not alter the
zero-temperature full-viability definition. The current
\code{--write-all-points} option does not bypass
Eq.~\eqref{eq:ewptgate}: the active caller still checks EWPT eligibility.
\src{G}\src{J}

\subsection{Qualitative baryogenesis and GW candidates}
The thermal coordinates are $(w_1,w_x,w_s)$. Symmetry-related sign copies
are mapped to positive representatives before jumps are evaluated, avoiding
spurious transitions between equivalent copies. An EW-entry pair requires
\begin{equation}
 |w_{1,\mathrm{false}}|<5\GeV\le|w_{1,\mathrm{true}}|.
\end{equation}
The field thresholds for classifying EW, $X$ and $S$ breaking are
$5$, $1$ and $1\GeV$, respectively, with exposed command-line overrides.
The candidate definitions use available first-order-transition pairs from
CalcTemps. The separate fully symmetric \code{SYM} phase label requires
all three field magnitudes below $1\GeV$; it is distinct from the
$|w_{1,\rm false}|<5\GeV$ EW-entry condition. \src{J}\src{A}

\begin{longtable}{@{}p{0.27\textwidth}p{0.67\textwidth}@{}}
\toprule
Flag & Positive criterion\\\midrule\endhead
\code{ewpt_baryo_candidate} & At least one EW-symmetric to EW-broken
FOPT at $T_c$ with
$\Delta v_{\rm EW}(T_c)/T_c>1$, where
$\Delta v_{\rm EW}=\bigl||w_{1,\rm true}|-|w_{1,\rm false}|\bigr|$.\\
\code{ewpt_gw_candidate} & At least one FOPT at an available
$T_c$, $T_n$ or $T_p$ with
$\Delta v_{\rm field}/T>1$, where
$\Delta v_{\rm field}=\sqrt{(\Delta|w_1|)^2+(\Delta|w_x|)^2+(\Delta|w_s|)^2}$.\\
\bottomrule
\end{longtable}
Both thresholds are \emph{strict}: equality at one fails. A larger
nucleation jump cannot rescue an insufficient critical EW-entry jump for
the baryogenesis flag. The GW flag can be positive from any of its three
temperatures. Percolation and completion are stored as diagnostics and do
not veto positive candidate evidence. These flags do not enter $V$.
\src{A}

A negative baryogenesis verdict requires adequate critical-temperature
assessment. A negative GW verdict also needs assessed nucleation and
percolation outcomes, since either could yield positive evidence.
The assessment requires successful NLO stability and tracing, and assessed
coexistence information; a successful no-coexistence-pair outcome can
establish a negative. For each relevant temperature, \code{not_met}
is an assessed outcome, while missing/failed calculations remain unknown.
Positive evidence is retained even if another part of the calculation is
incomplete. \src{A}

These are qualitative transition screens. They do not calculate a
sphaleron rate, a baryon asymmetry or a GW spectrum, and do not establish
that the Universe actually traverses a particular transition. The
finite-temperature potential and jump ratios retain their gauge and
approximation dependence.

\subsection{Equilibrium history and freeze-out consistency}
The separate phase analysis compares traced local minima at common
temperatures, merges symmetry copies and refines potential crossings.
The default target crossing bracket is at most $10^{-3}\GeV$ where
numerically resolvable. The resulting ordering is equilibrium ordering
among the traced basins, rather than proof that all minima were found or
that a cosmological transition completed. CalcTemps return codes,
MinimaTracer/refinement outcomes, NLO stability and phase ordering have
separate status fields. \src{A}\src{X}

The thermal dark-matter diagnostics include:
\begin{enumerate}
\item \code{dm_relic_z2_freezeout_compatible}: a necessary equilibrium
$Z_2$-compatibility screen around and after nominal freeze-out. Any sampled
$X$-broken state at $T\le T_f$ makes it false. A true verdict requires
coverage down to $T\simeq0$ (within $10^{-9}\GeV$), a resolved unbroken
state at $T_f$, and every sampled state at $T\le T_f$ unbroken. Otherwise
the verdict remains unavailable. Broken-interval endpoints are sampled
temperatures; restoration is given as a bracket. The helper consumes the
supplied lowest-potential branch, so this flag alone does not certify the
equilibrium reconstruction used in a historical input.
\item VEV variation over $[T_f/2,2T_f]$. The largest fractional changes
relative to the traced $T=0$ VEVs are recorded, and a change
$\ge10\%$ sets the VEV-shift diagnostic. Missing coverage and a phase
boundary crossing are treated explicitly; the code does not interpolate
between different minima to fabricate a smooth branch.
\item Fixed-zero-temperature-mixing coupling proxies,
\begin{align}
 K_{133}(T)&=\tfrac12[\lambda_{\Phi X}|w_1(T)|c-\lambda_{SX}|w_s(T)|s],\nonumber\\
 K_{233}(T)&=\tfrac12[\lambda_{\Phi X}|w_1(T)|s+\lambda_{SX}|w_s(T)|c].
\end{align}
These proxies require one resolved EW-broken ($|w_1|\ge5\GeV$),
$X$-unbroken ($|w_x|<1\GeV$) branch across the window and its
interpolation neighbours. Absolute changes, fractional changes and
cancellations are recorded separately. Writing $K_a$ for the two portal
terms, fractional changes are unavailable when
$|K(0)|\le10^{-12}\max(\sum_a|K_a(0)|,1\GeV)$.
Separately, a cancellation ratio
$|\sum_a K_a|/\sum_a|K_a|<0.1$, at zero temperature or a sampled window
point, flags cancellation sensitivity; that flag alone does not suppress
the fractional result.
\item Signed mass gaps $M_i-2M_3$ and the nonnegative diagnostics
$|M_i-2M_3|/T_f$ and $|M_i-2M_3|/\Gamma_i$, using backend mediator widths.
\end{enumerate}
These diagnostics do not alter the stored aggregate DM verdict, recompute
thermal mass eigenstates, or solve a temperature-dependent Boltzmann
evolution. \src{X}\src{K}

\section{Reading the plotting suite}
\label{sec:plots}
\subsection{Current and historical masks}
The main suite reads flags from its input TSV, including files with a
\code{.dat} suffix. It does not rerun the constraint calculations.
For rows labelled \code{trsm_constraints_v2}, theory is \code{thc}
and experimental is the stored \code{experimental_subset}. For legacy
rows, theory is \code{evo} $\land$ \code{thc} and experimental is
\code{hb} $\land$ \code{hs} $\land$ \code{ewpo} $\land$ \code{wmass}.
Full and non-DM viability then
include the stored flavour flag as in Eq.~\eqref{eq:viability}.
Therefore an old file without flavour assessments cannot acquire current
full viability merely by being replotted. \src{P}

The four-way DM/experimental maps classify the independent pair $(D,E)$.
Relic-only and DD-only maps instead use an available, nonexcluded component.
The main suite's detailed DM availability mask requires non-null entries
for \code{dm_relic_excluded}, \code{dm_direct_detection_excluded},
\code{dm_indirect_available} and \code{dm_indirect_detection_excluded}.
This bookkeeping mask is distinct from the aggregate DM verdict and from
availability of an actual line bound. CMB has additional dedicated flags.

\subsection{Independent v2 overlays and thermal comparisons}
\code{v2_candidate_subsets} shows baryogenesis and GW candidates, with
separate intersections with DM, global tree vacuum and RG subsets.
The RG overlay requires all of integration success, running BFB and running
unitarity. These three overlays are not successive cuts. Unknown candidate
and experimental states receive distinct markers.
\code{v2_jump_temperatures} compares critical and nucleation/percolation
jumps for the same transition, with missing partners counted.
\code{v2_freezeout_restoration} and \code{v2_thermal_sensitivity}
display the diagnostics in the preceding section. \src{V}

The older selected-transition quantity \code{ewpt_ew_true_over_T}
is a separate $v_{\rm true}/T$ diagnostic. Its temperature priority is
nucleation, percolation, completion, then critical. It must not be
substituted for the critical EW-entry jump in the baryogenesis definition.
Its implemented strong/weak split is $v/T>1$ versus $v/T\le1$.
\src{P}\src{G}

\subsection{Restrictions needed to display signal rates}
Cascade and MadGraph mono-Higgs/mono-$Z$ panels use either full viability
or non-DM viability, as stated in each plot, then require finite positive
rates for display. The YR4 invisible-$h_2$ signal panels require full
viability $V$, followed by physical finite inputs, $M_2>2M_3$, a positive
$\BR(h_2\to h_3h_3)$, supported mass interpolation and a positive
finite rate. The displayed narrow-width production proxy is
\begin{equation}
 \sigma_{\rm inv}^{p}(13.6\,\mathrm{TeV})
 =\sigma_{\rm SM}^{p}(M_2)k_2^2\BR(h_2\to h_3h_3),
 \qquad p\in\{\mathrm{ggF},\mathrm{VBF}\}.
\end{equation}
The table is not extrapolated. These requirements select where the
observable can be shown; they do not add a new exclusion to $V$.
Raw HL-LHC event proxies use $3000\,\mathrm{fb}^{-1}$ without acceptance
or efficiency factors. \src{P}

The width categories are
\begin{equation}
 \frac{\Gamma_2}{M_2}<0.01,\qquad
 0.01\le\frac{\Gamma_2}{M_2}<0.10,\qquad
 \frac{\Gamma_2}{M_2}\ge0.10.
\end{equation}
Broad points remain in the plots; the width class is not a veto. Physical
$\Gamma_i/M_i$ and $c\tau_i$ are recorded independently of the provider
width floor. No prompt-decay or upper-width cut is part of full viability.

Mass-plane lines such as $M_2=2M_3$, $M_3=2M_2$,
$M_3=M_2+125\GeV$ and $M_3=1000\GeV$ are guides, not cuts applied
by the plotter. In particular, the guide's nominal $125\GeV$ gap differs
from the generator's $M_1=125.09\GeV$, and $M_3=2M_2$ is not an
$h_3$ decay threshold in this stable-DM branch. Resonance-coloured panels
do not themselves impose a resonance window. The flavour zoom's
$4\le M_2/\mathrm{GeV}\le9.2$ restriction selects the displayed window.
\src{P}\src{F}

\subsection{Wording and legacy-code caveats}
Some retained labels predate the current masks. When interpreting a plot
or its summary, the following code definitions take precedence:
\begin{itemize}
\item The summary text may say theory is \code{evo} $\land$ \code{thc}; for v2
the actual mask is \code{thc}. Its generic experimental text also omits
the coverage-aware use of \code{experimental_subset}.
\item Some older EWPT bar labels say strong $v/T\ge1$ and weak $v/T<1$;
the implemented split is strictly strong $>1$ and weak $\le1$.
\item The main suite's BSMPT ``success'' grouping includes the statuses
\code{success}, \code{assessed} and \code{incomplete}. It does not
guarantee a fully assessed thermal history; use the detailed v2 statuses.
\item In the general $v_x\ne0$ evaluator \code{evaluate_trsm_point},
the theory and evolution calls are commented out and the code explicitly
sets \code{thc=evo=True}. The existence of the general theoretical
functions therefore does not mean that this historical branch imposes them.
The active $v_x=0$ implementation described here evaluates them as specified.
\item \code{plot_trsm_observables.py} applies finite-value and ordinary
log-axis positivity requirements, not a new physics selection.
Its Eq.~4.18 quantities are observables unless separately requested as
an EWPT preselection. The historical \code{plot_trsm_points.py} reads a
hardcoded merged input and its ``Viable points'' title does not define a
current viability predicate.
\end{itemize}
These are interpretive notes about existing code; this document does not
change any constraint, data file or plotting behaviour. \src{P}\src{O}

\section{Reproducing a selection and locating its implementation}
\label{sec:sources}
From the repository's configured Python environment, the comprehensive
suite can be run on a chosen scan as follows:
\begin{quote}\ttfamily
python plot\_trsm\_constraint\_suite.py INPUT.tsv --format both
\end{quote}
Its default destination is \code{plots/<input-stem>_constraints} inside
the repository. \code{index.html} and \code{constraint_summary.tsv}
provide the figure inventory and selection counts; \code{v2-index.html}
adds the diagnostic views when the input contains v2 rows.
The cumulative entries always use the stored
input rows as their denominator. A historical input may already have been
filtered before plotting, so those counts do not reconstruct absent draws.

For a scientific comparison, retain the input and metadata sidecar, physics
version, source and executable fingerprints, Higgs dataset revisions,
direct-detection table, CMB setting and actual sampling configuration.
The current defaults cannot retrospectively identify an older scan's
settings. Adding only flavour assessments with
\code{reevaluate_trsm_flavour.py} does not revalidate its other historical
verdicts. \src{R}\src{F}

\small
\begin{longtable}{@{}p{0.055\textwidth}p{0.40\textwidth}p{0.485\textwidth}@{}}
\toprule
Key & Repository-relative source & Relevant implementation\\
\midrule\endhead
\hypertarget{src:G}{G} & \code{generate_trsm_points.py} &
\code{_evaluate_trsm_point_vxzero} (1556--1654);
\code{run_ewpt_if_requested} (1312 onward);
\code{sample_random_masses}, \code{draw_random_vxzero_candidate};
range defaults (1667 onward).\\
\hypertarget{src:R}{R} & \code{trsm_constraint_profile.py} &
\code{profile_updates}: experimental and EWPT subsets, coverage;
\code{physics_manifest}: source/dataset/executable provenance.\\
\hypertarget{src:P}{P} & \code{plot_trsm_constraint_suite.py} &
\code{cumulative_constraint_masks} (2029--2049);
\code{load_scan} (2150 onward; masks at 2264--2279);
\code{constraint_bar_metrics}, \code{build_summary};
\code{derive_signal_observables}, \code{derive_bsmpt_results}.\\
\hypertarget{src:V}{V} & \code{plot_trsm_v2.py} &
\code{render_suite}: nullable candidate/subset overlays and
same-transition temperature comparisons.\\
\hypertarget{src:I}{I} & \code{trsm_inputs.py}; \code{config/sm-inputs-v2.json} &
Shared inputs, relic/abundance constants, RG boundary convention and
\code{nullable_and}.\\
\hypertarget{src:C}{C} & \code{config/next-scan-v2.json}; \code{tools/run_next_scan.py} &
Explicit campaign settings and construction of the EWPT-enabled command.\\
\hypertarget{src:T}{T} & \code{test_trsm_theory_constraints.py} &
\code{_copositive_quartic_matrix}, \code{_unitarity_eigenvalues},
\code{_passes_tree_level_constraints}, \code{theory_constraints_vxzero}.\\
\hypertarget{src:Q}{Q} & \code{trsm_theory_diagnostics.py}; \code{test_trsm_evolution.py} &
\code{potential_parameters}, \code{assess_vacuum},
\code{assess_running}, \code{theory_diagnostics}; one-loop \code{rhs}.\\
\hypertarget{src:W}{W} & \code{generate_trsm_info.py} &
\code{Rmatrix}, \code{vxzero_portal_couplings},
\code{scalar_to_identical_scalar_width}, \code{vxzero_invisible_decay_info}.\\
\hypertarget{src:H}{H} & \code{test_trsm_higgstools.py} &
\code{analyze_parampoint} (422--496), \code{_sm_reference_chi2},
\code{_set_base_decays}; physical BRs and provider width floor.\\
\hypertarget{src:E}{E} & \code{singlet_EWPO.py}; \code{datafiles/Tania_MW_SnowmassWhitepaper.dat} &
\code{check_EWPO_wU}, \code{calc_chisquared_correlated_wU}; fit inputs
(255--267); \code{check_wmass_tania} (46--72).\\
\hypertarget{src:F}{F} & \code{trsm_flavour.py}; \code{flavour/flavour_observables.py};
\code{flavour/flavour_inputs.py}; \code{plot_trsm_flavour.py} &
\code{assess_scalar}, \code{combine_flavour},
\code{EvaluateUpsLeptonBounds}, \code{BRUpsgammahiNorm};
digitized limits and flavour plot categories.\\
\hypertarget{src:D}{D} & \code{test_trsm_DM.py}; \code{DM/main.c} &
\code{DMSummary.passed} (210--218), \code{summarize_dm_result},
\code{assess_indirect_limit}, \code{test_dm}; solver, spectra and line output.\\
\hypertarget{src:L}{L} & \code{trsm_direct_detection.py};
\code{DM/data/lz2024/observed-si-v2.json} &
Default numerical DD table, validation, interpolation and units.\\
\hypertarget{src:B}{B} & \code{trsm_cmb.py}; \code{DM/main.c} &
\code{assess_cmb_limit}, configuration/capability checks;
\code{calcSpectrum} and \code{PlanckCMB} driver calls.\\
\hypertarget{src:J}{J} & \code{test_trsm_ewpt.py} &
\code{check_eq_4_18} (661--712), \code{EWPTConfig},
\code{run_trsm_ewpt}; extraction of thermal transition strengths.\\
\hypertarget{src:A}{A} & \code{ewpt_entry_criterion.py}; \code{ewpt_assessment.py} &
EW-entry/GW candidate definitions, strict thresholds, per-temperature
assessment and positive/negative/unknown logic.\\
\hypertarget{src:X}{X} & \code{ewpt_equilibrium.py}; \code{ewpt_x_history.py};
\code{BSMPT/src/prog/PhaseProbe.cpp} &
Common-temperature ordering, symmetry equivalence, crossing refinement,
$X$-broken intervals and freeze-out compatibility.\\
\hypertarget{src:K}{K} & \code{dm_thermal_relic_diagnostic.py} &
\code{resonance_proximity_updates}, \code{thermal_vev_updates},
\code{coupling_sensitivity_updates}.\\
\hypertarget{src:O}{O} & \code{plot_trsm_observables.py}; \code{plot_trsm_points.py} &
\code{rows_for_plot}: display validity; historical hardcoded-input plotter.\\
\bottomrule
\end{longtable}
\normalsize
Line numbers refer to the inspected working tree and may move in later
revisions. Files prefixed \code{test_trsm_} in this index include the
production physics routines imported by the generator; their names do not
mean that they are only test fixtures.

\end{document}
